\documentclass[addpoints, 12pt]{exam}
\usepackage[utf8]{inputenc}
\usepackage[spanish]{babel}
\usepackage{amsmath}
\usepackage{multicol}
\usepackage{graphicx}
\usepackage{booktabs}
\usepackage{tabularx}

% Configuración para líneas de rellenar
\setlength\fillinlinelength{3.5cm}

\begin{document}

\begin{questions}



\question Clasifica los siguientes fenómenos escribiendo una \textbf{F} si se trata de un cambio físico o una \textbf{Q} si es un cambio químico:
\begin{parts}
    \part \fillin[F] Se derrite un cubo de hielo.
    \part \fillin[Q] Digestión de alimentos en nuestro estómago.
    \part \fillin[F] Se rompe una hoja de papel en varios pedazos.
    \part \fillin[Q] Se derrite (y quema) una vela.
    \part \fillin[F] La lluvia.
\end{parts}

















\vspace{0.5cm}


\vspace{0.5cm}

\fullwidth{\Large \textbf{Parte II. Prefijos y Múltiplos}}






\vspace{0.3cm}

\question Clasifica los siguientes instrumentos de medición según la magnitud física para la que se emplean: \textit{Calibrador Vernier, Báscula, Termómetro infrarrojo, Cronómetro, Dinamómetro, Cinta métrica, Termómetro clínico, Reloj de sol.}

\begin{center}
\renewcommand{\arraystretch}{1.5}
\begin{tabularx}{0.9\textwidth}{X X X X}
\toprule
\textbf{Longitud} & \textbf{Masa} & \textbf{Temperatura} & \textbf{Tiempo / Fuerza} \\
\midrule
\vspace{2cm} & \vspace{2cm} & \vspace{2cm} & \vspace{2cm} \\
\bottomrule
\end{tabularx}
\end{center}
















\question Con base en la lectura siguiente, responde brevemente a las siguientes preguntas:

\begin{center}
    \fbox{\parbox{0.9\textwidth}{
    María Sklodowska. Mejor conocida como Marie Curie, fue la primera mujer en ganar un Premio Nobel y la primera persona en ganar dos premios Nobel en áreas científicas distintas. Nació el 7 de noviembre de 1867 en Varsovia. A los 15 años no pudo ingresar a la universidad de Varsovia porque no admitía mujeres, sin embargo, sus deseos de aprender no se apagaron e ingresó a la ``Universidad Volante'', una escuela clandestina para mujeres que constantemente cambiaba de sede para escapar de la persecución rusa. En 1891, con apoyo de su hermana, Marie se matriculó en la Universidad de París para estudiar Física, Química y Matemáticas. En 1894 conoció a Pierre Curie, físico que se dedicaba al estudio de la radioactividad. Entablaron una gran amistad y se casaron. Con su esposo estudió la radiación del uranio y descubrieron el polonio y el radio. Desafortunadamente, a consecuencia de la exposición prolongada a sustancias radioactivas, enfermó de anemia aplástica y falleció el 4 de julio de 1934. Su legado; una de sus hijas; Iréne Joliot Curie, siguió el ejemplo de sus padres, estudió Física y Química y recibió el Premio Nobel de Química en 1935. Gracias a Marie, la Universidad admitió a 45 mujeres entre 1906 y 1934, haciendo a un lado sus antiguas restricciones de género. Marie Curie junto a su hija Iréne trabajando en su laboratorio en París, en abril de 1927.
    }}
\end{center}

\begin{parts}
    \part ¿Cuál es la diferencia entre conocimiento empírico y conocimiento científico?
    \makeemptybox{3cm}
    
    \part ¿Cómo beneficiaron a la humanidad los descubrimientos de la familia Curie?
    \makeemptybox{3cm}
\end{parts}


\question Indica a qué científico pertenece cada aportación (J.J. Thomson, Ernest Rutherford, Niels Bohr, o Ninguno):
\begin{parts}
    \part Demostró que los rayos catódicos son partículas cargadas negativamente y que provienen de los átomos de la materia. (\textbf{J.J. Thomson})
    \part Concibió un modelo de átomo conocido como ``panqué de pasas'' donde el pan representaba la carga positiva y las ``pasas'', las cargas negativas. (\textbf{J.J. Thomson})
    \part Con base en sus experimentos con materiales radiactivos, demostró que la carga positiva no era una especie de nube que envuelve a los electrones. (\textbf{Ernest Rutherford})
    \part Aceptando parcialmente el modelo de Rutherford, aseguró que los electrones sólo giran en órbitas determinadas alrededor del núcleo. (\textbf{Niels Bohr})
    \part Formuló un modelo que permitió explicar la existencia de líneas oscuras en el espectro luminoso del hidrógeno, y demostró que éstas tenían su origen en el intercambio de energía con los electrones en el núcleo. (\textbf{Niels Bohr})
    \part Concibió un modelo atómico planetario en el que el Sol representaba al núcleo y los planetas representaban los electrones que orbitaban alrededor de éste. (\textbf{Ernest Rutherford})
    \part Explicó que, cuando regresan a órbitas menos energéticas, los electrones deben desprenderse de una cantidad fija de energía y que ésta se manifiesta como luz. (\textbf{Niels Bohr})
    \part Aseguró que en la naturaleza la energía no es continua, sino que se presenta en múltiplos de una cantidad fija conocida como cuanto. (\textbf{Niels Bohr})
    \part Demostró que el plutonio, el uranio, el radio y el polonio son elementos radiactivos, pues emitían partículas alfa, beta y gama. (\textbf{Ninguno})
\end{parts}

\question Completa las oraciones sobre los modelos atómicos colocando la frase correcta:
\begin{parts}
    \part El modelo atómico de J.J. Thomson es conocido como el \fillin[modelo atómico del ``panque de pasas''] y permitió \fillin[el descubrimiento del electrón].
    
    \part El modelo atómico de \fillin[Ernest Rutherford] es el modelo planetario del átomo que permitió \fillin[el descubrimiento del núcleo atómico] formado por protones y \fillin[neutrones].
    
    \part El modelo atómico de \fillin[Niels Bohr] es conocido como el \fillin[modelo cuántico del átomo] y permitió \fillin[describir el átomo de hidrógeno] y \fillin[explicar los espectros luminosos].
\end{parts}

\question A partir de las siguientes descripciones visuales, identifica de qué modelo atómico se trata:
\begin{parts}
    \part Un esquema que muestra a un electrón saltando entre diferentes órbitas circulares concéntricas y absorbiendo o emitiendo energía en forma de luz:
    \begin{oneparchoices}
        \choice J.J. Thomson \choice Ernest Rutherford \CorrectChoice Niels Bohr
    \end{oneparchoices}
    
    \part Una esfera naranja continua de carga positiva que tiene incrustadas varias partículas pequeñas con carga negativa en su interior:
    \begin{oneparchoices}
        \CorrectChoice J.J. Thomson \choice Ernest Rutherford \choice Niels Bohr
    \end{oneparchoices}
\end{parts}


\question Completa las siguientes afirmaciones con la palabra correcta:
\begin{parts}
    \part La Física es una \fillin[ciencia] que utiliza el método \fillin[científico] para obtener sus conocimientos.
    \part La Física estudia la \fillin[naturaleza] y los fenómenos que ocurren en ella, incluso su origen y su posible final.
    \part También estudia el Universo y sus componentes fundamentales, como los \fillin[átomos] y las \fillin[partículas] que los forman.
    \part Los cambios que experimentan los objetos y que no alteran su composición se denominan cambios \fillin[físicos].
    \part La Física estudia la relación entre \fillin[la materia] y la energía, la interacción entre objetos, el cambio producido por \fillin[una fuerza] y dónde y cuándo tienen lugar los fenómenos, es decir, \fillin[el espacio y el tiempo].
\end{parts}

\question Clasifica los siguientes fenómenos escribiendo una \textbf{F} si se trata de un cambio físico o una \textbf{Q} si es un cambio químico:
\begin{parts}
    \part \fillin[F] Se derrite un cubo de hielo.
    \part \fillin[Q] Digestión de alimentos en nuestro estómago.
    \part \fillin[F] Se rompe una hoja de papel en varios pedazos.
    \part \fillin[Q] Se derrite (y quema) una vela.
    \part \fillin[F] La lluvia.
\end{parts}

\question Relaciona los siguientes conceptos con la rama de la Física que les corresponde. Escribe la letra en el paréntesis.
\begin{multicols}{2}
    \begin{enumerate}
        \item[a)] Física clásica
        \item[b)] Física moderna
        \item[c)] Cinemática
        \item[d)] Dinámica
        \item[e)] Termodinámica
        \item[f)] Acústica
    \end{enumerate}
    \columnbreak
    \begin{parts}
        \part (\fillin[a][1cm]) Estudia fenómenos a velocidades mucho menores a la de la luz.
        \part (\fillin[b][1cm]) Estudia fenómenos a nivel atómico y el origen del Universo.
        \part (\fillin[c][1cm]) Describe el movimiento de los objetos.
        \part (\fillin[d][1cm]) Estudia la causa del movimiento de los objetos (fuerza).
        \part (\fillin[e][1cm]) Estudia el calor y la temperatura.
        \part (\fillin[f][1cm]) Rama del movimiento ondulatorio que estudia el sonido.
    \end{parts}
\end{multicols}

\question Clasifica las siguientes ramas escribiendo \textbf{C} si pertenecen a la Física Clásica o \textbf{M} si pertenecen a la Física Moderna:
\begin{multicols}{2}
\begin{parts}
    \part \fillin[C][1cm] Mecánica
    \part \fillin[C][1cm] Electromagnetismo
    \part \fillin[C][1cm] Termodinámica
    \part \fillin[M][1cm] Relatividad
    \part \fillin[M][1cm] Mecánica cuántica
\end{parts}
\end{multicols}

\question Determina a qué rama de la física corresponde cada uno de los siguientes escenarios o sistemas de estudio:
\begin{parts}
    \part Un termómetro midiendo altas temperaturas expuesto al sol:
    \begin{oneparchoices}
        \choice \textbf{Termodinámica} \choice Óptica \choice Acústica \choice Cosmología
    \end{oneparchoices}
    
    \part Antenas de telecomunicaciones emitiendo señales:
    \begin{oneparchoices}
        \choice Dinámica \choice \textbf{Electromagnetismo} \choice Cinemática \choice Óptica
    \end{oneparchoices}
    
    \part La observación de una nebulosa en el espacio profundo:
    \begin{oneparchoices}
        \choice Termodinámica \choice Dinámica \choice \textbf{Cosmología} \choice Acústica
    \end{oneparchoices}
    
    \part Uso de un microscopio en un laboratorio:
    \begin{oneparchoices}
        \choice Electromagnetismo \choice Cosmología \choice Acústica \choice \textbf{Óptica}
    \end{oneparchoices}
    
    \part Automóviles de Fórmula 1 compitiendo en una pista (estudio de su movimiento):
    \begin{oneparchoices}
        \choice Termodinámica \choice \textbf{Cinemática} \choice Óptica \choice Electromagnetismo
    \end{oneparchoices}
    
    \part Una orquesta sinfónica en pleno concierto:
    \begin{oneparchoices}
        \choice Dinámica \choice \textbf{Acústica} \choice Cosmología \choice Cinemática
    \end{oneparchoices}
    
    \part Una persona empujando un automóvil descompuesto (estudio de las fuerzas):
    \begin{oneparchoices}
        \choice Termodinámica \choice Acústica \choice Óptica \choice \textbf{Dinámica}
    \end{oneparchoices}
\end{parts}

\vspace{0.5cm}
\fullwidth{\Large \textbf{Parte III: Mediciones y Unidades Fundamentales}}

\question Utiliza las siguientes palabras para completar las afirmaciones: \textit{comparar, cuantificar, medición, precisión, referencia}.
\begin{parts}
    \part Medir significa \fillin[comparar] dos cantidades con un instrumento de \fillin[medición].
    \part Una unidad de medida es una \fillin[referencia] que permite estimar una propiedad que se puede \fillin[cuantificar] con cierta exactitud.
    \part La forma de medir en la antigüedad era muy diferente a como lo hacemos en la actualidad, la principal diferencia es la \fillin[precisión] del instrumento con el que se mide.
\end{parts}

\question Subraya o marca la(s) respuesta(s) correcta(s) para cada enunciado (puede haber más de una):
\begin{parts}
    \part Son algunas unidades de medida que se utilizaban en la antigüedad:
    \begin{checkboxes}
        \choice Libra
        \CorrectChoice Pie
        \choice Yarda
        \CorrectChoice Codo
        \choice Metro
    \end{checkboxes}
    
    \part Los mexicas tenían unidades de medida de longitud, como el \rule{2cm}{0.15mm} que era un palo de madera que medía aproximadamente dos metros y medio.
    \begin{checkboxes}
        \choice testal
        \choice cenequeztzalli
        \choice omitl
        \CorrectChoice tlalcuahuitl
    \end{checkboxes}
    
    \part Son medidas de longitud usadas en México en la época colonial:
    \begin{checkboxes}
        \CorrectChoice Legua
        \choice Testal
        \CorrectChoice Vara
        \choice Arroba
    \end{checkboxes}
    
    \part La \rule{2cm}{0.15mm} es la medida de la palma de la mano extendida.
    \begin{checkboxes}
        \choice legua
        \choice vara
        \CorrectChoice cuarta
        \choice yarda
    \end{checkboxes}
    
    \part La arroba fue una unidad de medida de \rule{2cm}{0.15mm} en la época colonial.
    \begin{checkboxes}
        \choice moral
        \CorrectChoice peso
        \choice técnico
        \choice práctico
    \end{checkboxes}
\end{parts}

\question Clasifica las siguientes unidades de medida históricas según la magnitud física que evaluaban (Capacidad, Longitud o Peso):
\begin{itemize}
    \item Barril: \fillin[Capacidad][3cm]
    \item Codo: \fillin[Longitud][3cm]
    \item Arroba: \fillin[Peso][3cm]
    \item Pinta: \fillin[Capacidad][3cm]
    \item Vara: \fillin[Longitud][3cm]
    \item Legua: \fillin[Longitud][3cm]
    \item Botija: \fillin[Capacidad][3cm]
    \item Testal: \fillin[Peso][3cm]
\end{itemize}


\end{questions}
\end{document}